\documentclass[10pt,a4paper]{amsart}
\usepackage{amsmath,amssymb,amsthm,graphicx}
\usepackage{verbatim}
\newenvironment{reference}{\comment}{\endcomment}
%\newenvironment{reference}{}{}
\newenvironment{slogan}{\comment}{\endcomment}
\newenvironment{history}{\comment}{\endcomment}
% For commutative diagrams we use Xy-pic
\usepackage[all]{xy}

% We use 2cell for 2-commutative diagrams.
\xyoption{2cell}
% Theorem environments.
%
\theoremstyle{plain}
\newtheorem{theorem}[subsection]{Theorem}
\newtheorem{proposition}[subsection]{Proposition}
\newtheorem{lemma}[subsection]{Lemma}

\theoremstyle{definition}
\newtheorem{definition}[subsection]{Definition}
\newtheorem{example}[subsection]{Example}
\newtheorem{exercise}[subsection]{Exercise}
\newtheorem{situation}[subsection]{Situation}

\theoremstyle{remark}
\newtheorem{remark}[subsection]{Remark}
\newtheorem{remarks}[subsection]{Remarks}

\numberwithin{equation}{subsection}

% Macros
%
\def\lim{\mathop{\mathrm{lim}}\nolimits}
\def\colim{\mathop{\mathrm{colim}}\nolimits}
\def\Spec{\mathop{\mathrm{Spec}}}
\def\Hom{\mathop{\mathrm{Hom}}\nolimits}
\def\Ext{\mathop{\mathrm{Ext}}\nolimits}
\def\SheafHom{\mathop{\mathcal{H}\!\mathit{om}}\nolimits}
\def\SheafExt{\mathop{\mathcal{E}\!\mathit{xt}}\nolimits}
\def\Sch{\mathit{Sch}}
\def\Mor{\mathop{\mathrm{Mor}}\nolimits}
\def\Ob{\mathop{\mathrm{Ob}}\nolimits}
\def\Sh{\mathop{\mathit{Sh}}\nolimits}
\def\NL{\mathop{N\!L}\nolimits}
\def\CH{\mathop{\mathrm{CH}}\nolimits}
\def\proetale{{pro\text{-}\acute{e}tale}}
\def\etale{{\acute{e}tale}}
\def\QCoh{\mathit{QCoh}}
\def\Ker{\mathop{\mathrm{Ker}}}
\def\Im{\mathop{\mathrm{Im}}}
\def\Coker{\mathop{\mathrm{Coker}}}
\def\Coim{\mathop{\mathrm{Coim}}}

% Boxtimes
%
\DeclareMathSymbol{\boxtimes}{\mathbin}{AMSa}{"02}

%
% Macros for moduli stacks/spaces
%
\def\QCohstack{\mathcal{QC}\!\mathit{oh}}
\def\Cohstack{\mathcal{C}\!\mathit{oh}}
\def\Spacesstack{\mathcal{S}\!\mathit{paces}}
\def\Quotfunctor{\mathrm{Quot}}
\def\Hilbfunctor{\mathrm{Hilb}}
\def\Curvesstack{\mathcal{C}\!\mathit{urves}}
\def\Polarizedstack{\mathcal{P}\!\mathit{olarized}}
\def\Complexesstack{\mathcal{C}\!\mathit{omplexes}}
% \Pic is the operator that assigns to X its picard group, usage \Pic(X)
% \Picardstack_{X/B} denotes the Picard stack of X over B
% \Picardfunctor_{X/B} denotes the Picard functor of X over B
\def\Pic{\mathop{\mathrm{Pic}}\nolimits}
\def\Picardstack{\mathcal{P}\!\mathit{ic}}
\def\Picardfunctor{\mathrm{Pic}}
\def\Deformationcategory{\mathcal{D}\!\mathit{ef}}
\usepackage{cleveref}
\crefname{originalreference}{original reference}{original references}
\Crefname{originalreference}{Original reference}{Original references}
\makeatletter
\newlabel{spaces-morphisms-lemma-unramified-flat-lfp-etale}{{spaces-morphisms-lemma-unramified-flat-lfp-etale}{source}{}{sourcepointer-1}{}}
\newlabel{spaces-morphisms-lemma-unramified-flat-lfp-etale@cref}{{[originalreference][0][]spaces-morphisms-lemma-unramified-flat-lfp-etale}{[1][0][]source}{}{}{}}
\makeatother
\crefname{originalreference}{original reference}{original references}
\Crefname{originalreference}{Original reference}{Original references}
\makeatletter
\newlabel{stacks-properties-item-etale}{{stacks-properties-item-etale}{source}{}{sourcepointer-2}{}}
\newlabel{stacks-properties-item-etale@cref}{{[originalreference][0][]stacks-properties-item-etale}{[1][0][]source}{}{}{}}
\makeatother
\crefname{originalreference}{original reference}{original references}
\Crefname{originalreference}{Original reference}{Original references}
\makeatletter
\newlabel{stacks-properties-remark-local-source-apply}{{stacks-properties-remark-local-source-apply}{source}{}{sourcepointer-3}{}}
\newlabel{stacks-properties-remark-local-source-apply@cref}{{[originalreference][0][]stacks-properties-remark-local-source-apply}{[1][0][]source}{}{}{}}
\makeatother
\title{A flat morphism of finite presentation which is unramified at a point of a field fibre is etale there}
\author{The Stacks Project Authors}
\date{stacks-project:stacks-morphisms.tex:local-snapshot-20261012; GNU FDL 1.2 or later}
\begin{document}
\maketitle
\noindent\textit{Editorial note.} All delivered mathematical text is the author's own, copied from the pinned Stacks Project chapter source sources/stacks-algebraic-stacks-morphisms/pinned-original-source.tex; no statement or proof step was rewritten, reordered or removed.\par\medskip
\noindent\textit{Editorial note.} Preamble provenance: the Stacks Project keeps one single shared preamble for every chapter; this card takes it directly from the staged sources/stacks-algebraic-stacks-morphisms/preamble.tex as three verbatim spans (lines 8-12, 14-18 and 154-217), so no derived file was written and nothing was added or redefined. The wrapper class and the base packages are supplied by the assembler.\par\medskip
\noindent\textit{Editorial note.} Bibliography provenance: neither the counted statement nor its proof contains a citation (this lane rescans the whole span before writing the card), so the card carries no bibliography and no entry was reconstructed.\par\medskip
\noindent\textit{Editorial note.} Selection provenance: this card was opened under the file-based control contract (CONTROL.md), with the mechanical pre-card gate run on its statement and proof interval before the card was written.\par\medskip
\noindent\textit{Allowed standard prerequisites (editorial).}\begin{itemize}
\item The existence and base change behaviour of the maximal open subspace on which a flat finitely presented morphism is etale. Invoked at the start of the counted proof (source lines 4567-4575) as cross-chapter pointers. Reduces the assertion to the membership of the given point in that open subspace.
\item The statement that a flat morphism of finite presentation which is unramified is etale. Invoked at the end of the counted proof (source lines 4576-4582) as a cross-chapter pointer. Gives etaleness of the base-changed morphism at the point.
\end{itemize}\medskip
\ifdefined\hypertarget\hypertarget{sourcepointer-1}{}\fi
\noindent\textit{Source reference.} Original reference spaces-morphisms-lemma-unramified-flat-lfp-etale: sources/stacks-algebraic-stacks-morphisms/pinned-original-source.tex, source lines 1131--1131. Navigation pointer only: the counted text refers to the Stacks item spaces-morphisms-lemma-unramified-flat-lfp-etale, which belongs to another chapter of the project and is not part of this card. The pointer carries the target's own label as its number because this chapter's staged snapshot does not contain the other chapter's tag table; its content is not reproduced here.\par\medskip
\ifdefined\hypertarget\hypertarget{sourcepointer-2}{}\fi
\noindent\textit{Source reference.} Original reference stacks-properties-item-etale: sources/stacks-algebraic-stacks-morphisms/pinned-original-source.tex, source lines 4572--4572. Navigation pointer only: the counted text refers to the Stacks item stacks-properties-item-etale, which belongs to another chapter of the project and is not part of this card. The pointer carries the target's own label as its number because this chapter's staged snapshot does not contain the other chapter's tag table; its content is not reproduced here.\par\medskip
\ifdefined\hypertarget\hypertarget{sourcepointer-3}{}\fi
\noindent\textit{Source reference.} Original reference stacks-properties-remark-local-source-apply: sources/stacks-algebraic-stacks-morphisms/pinned-original-source.tex, source lines 4270--4270. Navigation pointer only: the counted text refers to the Stacks item stacks-properties-remark-local-source-apply, which belongs to another chapter of the project and is not part of this card. The pointer carries the target's own label as its number because this chapter's staged snapshot does not contain the other chapter's tag table; its content is not reproduced here.\par\medskip
\begin{lemma}
\label{lemma-etale-at-point}
Let $\mathcal{X}$ be an algebraic stack. Consider a cartesian diagram
$$
\xymatrix{
U \ar[d] & F \ar[l]^p \ar[d] \\
\mathcal{X} & \Spec(k) \ar[l]
}
$$
where $U$ is an algebraic space, $k$ is a field, and $U \to \mathcal{X}$
is flat and locally of finite presentation. Let $z \in |F|$ be such that
$F \to \Spec(k)$ is unramified at $z$. Then, after replacing $U$ by
an open subspace containing $p(z)$, the morphism
$$
U \longrightarrow \mathcal{X}
$$
is \'etale.
\end{lemma}

\begin{proof}
Since $f : U \to \mathcal{X}$ is flat and locally of finite presentation
there exists a maximal open $W(f) \subset U$ such that the restriction
$f|_{W(f)} : W(f) \to \mathcal{X}$ is \'etale, see
Properties of Stacks, Remark
\ref{stacks-properties-remark-local-source-apply}
(\ref{stacks-properties-item-etale}).
Hence all we need to do is prove that $p(z)$ is a point of $W(f)$.
Moreover, the remark referenced above also shows the formation of $W(f)$
commutes with arbitrary base change by a morphism which is representable
by algebraic spaces. Hence it suffices to show that the morphism
$F \to \Spec(k)$ is \'etale at $z$. Since it is flat and locally
of finite presentation as a base change of $U \to \mathcal{X}$ and since
$F \to \Spec(k)$ is unramified at $z$ by assumption, this follows
from
Morphisms of Spaces,
Lemma \ref{spaces-morphisms-lemma-unramified-flat-lfp-etale}.
\end{proof}

\end{document}
